\documentclass[11pt]{amsart}

\usepackage{amsmath, amssymb, amsthm}
\usepackage{mathtools}
\usepackage{thmtools}
\usepackage{enumitem}
\usepackage{tikz-cd}
\usepackage{hyperref}
\usepackage[nameinlink]{cleveref}

% % Theorem environments
\declaretheorem[numberwithin=section]{theorem}
\declaretheorem[style=plain, sibling=theorem]{lemma, proposition, corollary}
\declaretheorem[style=definition]{definition, example, condition}
\declaretheorem[style=remark]{remark, note, observation, todo}


\title{April 19, 2026}

\begin{document}
\maketitle

Today I strated the day with the yeterday's unsolved burden on the bed; a relation between 3-cocycles on tetrahydra and tensor coefficients on 3-valent graph.
After having lunch, and watching Youtube little bit, I came to the school and got ready to study.
While I chose a playlist for today's research, I found a video saying that lying on the bed might be harmful to your brain cells.
Here, I record the summary of that video:

When you lie on the bed, more than 3 days, then the density of C-Reactive Protine increases in your body.
It implies that the secration of dophamin.
Hence, without the physical activity, there is a negative feedback making your body patient.
In this sense, the behavior of tending to lie on the bed is called the \emph{bed rotting}.

The way to overcome is suggested as follows:
\begin{enumerate}
    \item Grounding: To focus on your body senses, move your toe, feel your back and the temperatur of your room.
    It turns off the default mode network, which is the region of your brain activated when you are blank or you have random thoughts.
    Instead, it turns on the salience newtork.
    Your body always keep the balance between the activity of those networks.
    Hence, focusing on your senses activate the salience network and deactivate the default mode network.

    \item Dophamin switching: change the contents generating dophamine.
    No short form, but yes shopping.
    Turn contents you watcing (passive) to contents you are ingaging (active).

    \item $45^\circ$ rule and RAS stimulation: In your brainstem, there is a region called the reticular activating system(RAS).
    It is a kind of a switch turning on/off your state of awakening.
    Just wakening your torso to $45^\circ$ turns on the RAS.
    
    \item Environmental nudge: Change your environment.
    Turning light really helps you since your braind is sensitive to the light.
\end{enumerate}

More easial drill is introduced, \emph{micro action drill}.
\begin{enumerate}
    \item Audio cue: Listen wake-up songs; instruments with no lyric, fast tempo with larger volume.
    Lyrics stimulate your language center.
    
    \item 5-4-3-2-1: Count down 5 secs.
    At the zero, just kick out your bedding, not wake-up.
    The activity of kicking will stimulate your brain.

    \item Zombie walking: After kicking the bedding, just walk around a minute without any thinking and any purpose.
    When the exprience that 'Oh, it is not hard to wake-up' is accumulated, then it will be easier to wake-up.
    
\end{enumerate}

\par\vspace{1em}

I got back to the paper of Kashaev.
I first reviewed \cite{Kashaev-1994-Qdilog6j} with the following understanding.

Let $BU_q$ be the Borel part of $U_q(\mathfrak{sl}_2)$ generated $D$ and $E$ such that
\[
    \Delta(E) = E \otimes E, \quad \Delta(D) = 1 \otimes D + D \otimes E.
\]  
When $q$ is a root of unity of order $N$, irreducible modules of $BU_q$ is classified by the values of $E^N$ and $D^N$ which are cenral in $BU_q$.
Let consider only irrep $V_p$ where $E^N$ acts by 1, and $D^N$ acts by for some $z_p \in \mathbb{C}$.
Then, by the Clebsch--Gordan decomposition for $BU_q$, we have
\[
    V_p \otimes V_q \cong N \cdot V_{pq}
\]
with $z_{pq} = z_p + z_q$.
Moreover, $V_p \otimes V_q \otimes V_r$ has $N^2$ terms of $V_{pqr}$ and there are two ways for expanding it:
\[
    (V_p \otimes V_q) \otimes V_r \cong V_p \otimes V_q \otimes V_r \cong V_p \otimes (V_q \otimes V_r).
\]
Hence, the isomorphism of $BU_q$ modules between $N^2 \cdot V_{pqr}$ gives a $R$-matrix which satisfies the pentagon equation.

Denote this solution by $R(p, q, r)$.
The author claims in Section 3 that, by proper normalization, one can find an oeprator $\Psi_{p, q, r}$ satisfies the quantum dilogarithm identity:
\[
    \Psi_{p,q,r}(U) \Psi_{p,qr,s}(-UV) \Psi_{q,r,s}(V) = \Psi_{pq,r,s}(V) \Psi_{p,q,rs}(U) 
\]
where $U$ and $V$ are proper operators (see the reference).
Hence, the quantum dilogarithm is a suitable normalization of $R$.

With another normalization of $R$ gives an assignment on tetrahedra with vertex, edge, and face labels inducing the state-sum invariant of link complements.
The label of vertices is essentially a coloring of representations on edges.
Specifically, when two vertices are colored by $z_p$ and $z_q$, then edge connecting them can be regarded as colored by the representation $V_r$ with $z_r = \pm (z_p - z_q)$.
It is equivalent to consider $3$-valent graphs on $S^2$ with a region colored by $z_p$, and each edge is colored by $V_r$ when its adjacent area is colored by $z_p$ and $z_q$.

Getting back to the polyhedra side, the label of edges is used to normalized $R$, and the label of faces is used to indicating specific basis of the linear space where $R$ acting.

\par\vspace{1em}

After then I read \cite{Kashaev-1997-VC} and \cite{Kashaev-1995-LinkQudilog}.
The former explains the validity of the volume conjecture by following reasons:
\begin{itemize}
    \item There exists a physical theroy, Euclidean quantum 2+1 gravity with a negative cosmological constant, gives a topological invariant of knot complement whose classical limit reproduces the hyperbolic volume of the complement.
    
    \item The hyperbolic volume of ideal tetrahedra in 3d hyperbolic space can be expressed in terms of (the imaginary part of) the Euler's dilogarithm.
    Hence, the generalization of the dilogarithm, the quantum dilogarithm, should generalize the notion of the hyperbolic volume.
    
    \item The link invariant is constructed from the cyclic quantum dilogarithm by Kashaev.
    It is, in fact, a quantum generalization of the hyperbolic volume (shown in this paper).
    Especially, it is shown that the classical limit of link invariants, for $4_1$ and $6_1$, reproduces their hyperbolic volume in this paper.
\end{itemize}

In \cite{Kashaev-1995-LinkQudilog}, the author shows that the relation between the $R$-matrix solution of the pentagon relation and the $R$-matrix solution of the YBE.
Moreprecisely, $RT$-invariant constructed from the solution of YBE is equal to the state-sum invariant constructed from the solution of the pentagon relation.
The key is that the state-sum of the ball minus a tangle, which has cell complex of an octahedron, representing a crossing can be expressed into the solution of YBE.
This observation, the (colored) solution of the YBE factors out of 4 terms of the (colored) solution of the pentagon relation seems that a generalization of the $R$-matrix of the quantum double of a Hopf algebra factors to 4 terms of the canonical element of the Hensenberg double of it.

I finished today's research by following remarks.
\begin{itemize}
    \item the cyclic dilogarithm can be obtained by the limit of the quantum dilogarithm as $q$ to a root of unity.
    
    \item In \cite{MM-2001-cJandKashaev}, the authors showed that the $R$-matrix of Kasahev (solution of the YBE) is equivalent to the universal $R$-matrix of $U_q(\mathfrak{sl}_2)$ on its representation via the Fourier transformation.
    
    \item Above two facts, with the fact that (Kashaev mentioned) the quantum dilogarithm is precisely the canonical element in the Heisenberg double of $BU_q$, imply that the $R$-matrix on the Verma module with the Fourier basis gives the Kashaev invariant.
    
    \item However, there is some ambiguity that (as McPhail-Snyder commented) the colored Jones polynomial is obtained from trivally colored $R$-matrix while the Kashaev's invariant require a coloring on regions, edges, and vertices.
    It seems that the process of summing up might be relavant to this ambiguity.
\end{itemize}

\bibliographystyle{amsalpha}
% \nocite{*}
\bibliography{bibliography}
\end{document}
